\documentclass[11pt]{article}
\usepackage[T1]{fontenc}
\usepackage{lmodern}
\usepackage[margin=1in]{geometry}
\usepackage{amsmath,amssymb,amsthm,mathtools,mathrsfs,booktabs,microtype}
\usepackage[colorlinks=true,linkcolor=blue,citecolor=blue,urlcolor=blue]{hyperref}
\newcommand{\R}{\mathbb R}
\newcommand{\Q}{\mathbb Q}
\newcommand{\Z}{\mathbb Z}
\newcommand{\C}{\mathbb C}
\newcommand{\N}{\mathbb N}
\newcommand{\OO}{\mathcal O}
\newcommand{\wt}[1]{\sum_{i=1}^m w_i #1_i}
\newtheorem{theorem}{Theorem}[section]
\newtheorem{lemma}[theorem]{Lemma}
\newtheorem{proposition}[theorem]{Proposition}
\newtheorem{corollary}[theorem]{Corollary}
\theoremstyle{remark}
\newtheorem{remark}[theorem]{Remark}
\title{The irrationality exponent of\texorpdfstring{\\$\log(3/2)+\sqrt2$}{ log(3/2) + sqrt(2)} is at most four}
\author{Ryan Matthew Casper}
\date{October 12, 2026}

\begin{document}
\maketitle
\begin{abstract}
We extend the paired-embedding affine determinant argument to fixed
nonintegral algebraic data. A column factor clears the additional
denominators while preserving the original matrix and its selected minor.
The resulting logarithmic cost is absorbed by the dimension and
approximation-weight choices, before the matrix height grows. The
specialization $a=3/2$, $u=1$, $v=\sqrt2$ gives an irrationality exponent
at most four for $\log(3/2)+\sqrt2$. The extension uses the interpolation
and analytic estimates of the preceding affine development. We distinguish
the mathematical argument, the accompanying Lean declarations, and their
verification status; no assertion of optimality or historical priority
is made.
\end{abstract}

\section{Statement and contribution}
For an irrational real number $x$, its irrationality exponent $\mu(x)$
is the supremum of the positive real $\nu$ for which
$0<|x-p/q|<q^{-\nu}$ holds for infinitely many reduced fractions
$p/q$ with $q>0$. Our result is stated first as a literal uniform bound.

\begin{theorem}\label{thm:pilot}
For every real $\nu>4$, there is an integer $Q\ge2$ such that
\begin{equation}\label{eq:pilot}
 q^{-\nu}\le \left|\log(3/2)+\sqrt2-\frac pq\right|
 \qquad(p\in\Z,\ q\in\N,\ q\ge Q).
\end{equation}
Consequently $\log(3/2)+\sqrt2$ is irrational and
$\mu(\log(3/2)+\sqrt2)\le4$.
\end{theorem}

The threshold precedes both $p$ and $q$; no coprimality assumption is
needed in \eqref{eq:pilot}. The theorem does not identify the exponent
exactly, compute $Q$, or assert that all exponents equal to four are
excluded.

The preceding affine paper \cite{affine} treated algebraic-integer
data in a degree-four field with two compatible embeddings. Its
geometric and analytic statements do not require integrality. This note
replaces the arithmetic step, accounts for its cost in the same
parameter construction, and specializes the result to a rational base
that was excluded by the old arithmetic hypothesis. The underlying
interpolation and determinant estimates are inherited from the
OpenAI library, the rational- and quadratic-logarithm extensions, and
their affine adaptation
\cite{openai,affine}; they are not new results of this note.

One checked published comparison comes from Marcovecchio
\cite{marcovecchio}, whose nonquadraticity table gives
$\mu_2(\log(3/2))<9.460812$. With that paper's
minimal-polynomial height, this implies the rational-exponent bound
$18.921624$ for our shifted target; Section~\ref{sec:benchmark}
gives the conversion. This is an applicable benchmark, not a
certification of the best published bound for this exact number.

\section{The denominator-aware affine theorem}\label{sec:common}
Write $\operatorname{Emb}(F)=\operatorname{Hom}_{\Q\text{-alg}}(F,\C)$.
The ring of integers is denoted by $\OO_F$.

\begin{theorem}\label{thm:general}
Let $[F:\Q]=4$ and $a,u,v\in F$. Suppose an embedding $\sigma_0$
satisfies $\sigma_0(a)=b_0\in\R$, with $b_0>0$ and $b_0\ne1$.
Let $x\in\R\setminus\{0\}$. Suppose a selected set
$G\subset\operatorname{Emb}(F)$ has cardinality two and
\begin{equation}\label{eq:compatible}
 \exp\bigl(\sigma(u)(x-\sigma(v))\bigr)=\sigma(a),
 \qquad |\sigma(u)|\le2 \quad(\sigma\in G).
\end{equation}
Then, for every real $\nu>4$, there is $Q\in\N$, $Q\ge2$, such that
$q^{-\nu}\le |x-p/q|$ for every $p\in\Z$ and $q\in\N$, $q\ge Q$.
In particular, $x$ is irrational and $\mu(x)\le4$.
\end{theorem}

Other embeddings may also satisfy \eqref{eq:compatible}, and
$\sigma_0$ need not belong to $G$. Embeddings are counted as distinct
maps even when their values on $a,u,v$ coincide. The result does not
require the remaining two embeddings to be real.

Choose positive integers $d_a,d_n$ such that
\begin{equation}\label{eq:fixedclearing}
 d_a a,\quad d_n u,\quad d_n uv\ \in\OO_F.
\end{equation}
Every element of a number field has a positive integer multiple in its
ring of integers. Taking a product of denominators clears the last two
elements simultaneously. These choices depend only on the fixed
algebraic data, not on any approximating fractions or matrix height.

\subsection{The inherited analytic and geometric input}
We state the precise part of the existing argument that is reused.
For a fixed $\nu>4$, suppose there are approximations
$|x-p/q|<q^{-\nu}$ with arbitrarily large $q$. Choose finitely many
$p_i/q_i$ with $q_i\ge2$ and $p_i\ne0$, and put
$w_i=\lceil\log q_i\rceil$, $w_* = \min_i w_i>0$.
Nonzero numerators are available by increasing the initial threshold,
since $x\ne0$.
There are rational shape parameters
\begin{equation}\label{eq:shape}
0<\theta<A<B<1,\quad C>1,\quad B<C^{-1},\quad B<C\theta<1,
\end{equation}
and $0<\eta<1$, for which
\begin{equation}\label{eq:gap}
g=\nu\bigl(A(1-\eta)-\theta\bigr)-2(1-\theta)>0.
\end{equation}
Choose $F_0>2/\theta$, a dimension $m\ge1$, and set
\begin{equation}\label{eq:dimension}
K=\lfloor C^m\rfloor,\qquad w_0=B^{-m},\qquad
v_0=2K\theta^m w_0.
\end{equation}
The finitely many weights $w_i$ are then chosen successively large
enough to satisfy the weight-separation hypotheses of weighted
logarithmic interpolation in \cite[Sections 4--5]{affine}. These
hypotheses compare finitely many products of earlier weights; an
arbitrarily large next denominator suffices. The strict volume ratio is
$Kw_0\theta^m/v_0=1/2$.

For $H>0$, use columns $(h,\gamma)\in\N\times\N^m$ with
$w_0h+\wt\gamma\le H$, and rows $(j,s,\beta)$ with
$0\le j<K$ and $v_0s+\wt\beta/\theta<H$. The number of rows is $M=M(H)>0$.
Define
\begin{align}
 T_i&=\left\lceil F_0w_i/v_0\right\rceil,&
 L_i(t)&=\sum_{1\le k<T_i}\frac{(-1)^{k+1}}k t^k,\label{eq:truncation}\\
 r_{\sigma i}&=\sigma(u)(p_i/q_i-\sigma(v)),&
 z_\sigma&=\sigma(a).\notag
\end{align}
The matrix is
\begin{equation}\label{eq:matrix}
\mathcal M_\sigma(H)_{(j,s,\beta),(h,\gamma)}
=z_\sigma^{jh}[t^sZ^\beta](1+t)^h
  \prod_{i=1}^m(Z_i+jr_{\sigma i}+L_i(t))^{\gamma_i}.
\end{equation}
Weighted interpolation gives full row rank at $\sigma_0$ for a cofinal
set of heights. Choose one nonzero full-row minor there and use its
\emph{same columns} at every embedding; its determinant is $\Delta_\sigma$.
Write
\[
b_H=\frac1{MH}\sum_{\text{rows }r}\sum_iw_i\beta_{ri},
\qquad 0\le b_H\le\theta.
\]

Fix bounds $\kappa,R\ge0$ and $\rho\ge100$ such that, at all embeddings,
\begin{align*}
\log\max(1,|\sigma(a)|)&\le\kappa,\\
|\sigma(u)|(|x|+1+|\sigma(v)|)&\le R,\\
100|x|&\le\rho,\qquad
100|\sigma(u)(x-\sigma(v))|\le\rho.
\end{align*}
Finite sums or maxima provide such bounds before the approximations
are selected. Put $\Lambda=\log4+4$ and $c=\log2/4$, and define
\begin{align}
 E_0&=\frac{\Lambda F_0m}{v_0}+\Lambda\sum_i\frac1{w_i}
                 +\frac\theta{w_*},\label{eq:costs}\\
 U(S)&=\frac{\log2}{v_0}+\frac{\kappa K+\log(3/2)}{w_0}
             +\frac{\log(2(KS+2))}{w_*},\notag\\
 V&=\frac\nu{F_0}+\frac{2\log2}{v_0}+\frac{\rho K}{w_0}
 +\frac{\log4+\log(2K)+\nu+\log(200K)+\log(\rho/100)}{w_*}.\notag
\end{align}
The inherited estimates \cite[Sections 7--8]{affine} assert that
there are remainders $r_a(H),r_c(H)\to0$ and a collision rate
$C_H\to C_\infty$, where
\begin{equation}\label{eq:collision}
C_\infty=\frac{c\eta^2K\theta^m}{(m+1)v_0A^m}
        =\frac{c\eta^2}{2(m+1)}(B/A)^m,
\end{equation}
such that every nonzero selected minor satisfies
\begin{align}
\frac{\log|\Delta_\sigma|}{MH}
&\le V+r_a(H)+\frac{\log2}{w_*}
 +\max\{-C_H,-\nu(A(1-\eta)-b_H)\},&&\sigma\in G,\label{eq:good}\\
\frac{\log|\Delta_\sigma|}{MH}
&\le U(R)+r_c(H),&&\sigma\notin G.\label{eq:coarse}
\end{align}
These are estimates for the original matrix \eqref{eq:matrix}; no
clearing factors are built into them. Their hypotheses are bounds on
complex bases and centres, the exponential identities, the approximation
errors and the numerical parameters. Algebraic integrality is absent.
Likewise, interpolation permits arbitrary additive centres. This is
why these statements remain applicable to the present data.

\section{Clearing denominators without changing the minor}
Let $D_i=\operatorname{lcm}(1,\ldots,T_i)$ and
$e_i=\lfloor H/w_i\rfloor$. The original row and column factors are
\[
\prod_iD_i^{e_i}q_i^{-\beta_i}
\quad\hbox{and}\quad \prod_iq_i^{\gamma_i}.
\]
Their product over the selected square matrix is a positive rational
number $\mathscr C_H$. The elementary clearing-cost calculation,
which is an inequality about these numerical factors, gives
\begin{equation}\label{eq:oldcost}
\log\mathscr C_H\le MH(1-b_H+E_0).
\end{equation}
It does not itself assert integrality for the present data.

Add to each selected column $(h,\gamma)$ the integer factor
\begin{equation}\label{eq:extracolumn}
c(h,\gamma)=d_a^{Kh}d_n^{\sum_i\gamma_i},
\end{equation}
and let $A_H$ be the product of these extra factors over the selected
columns. The exponent $Kh$ is a convenient upper bound; replacing it
by $(K-1)h$ is possible but unnecessary.

\begin{lemma}\label{lem:lift}
There is a single nonzero algebraic integer $Z_H\in\OO_F$ such that
\begin{equation}\label{eq:lift}
\sigma(Z_H)=\mathscr C_H A_H\Delta_\sigma
\qquad(\sigma\in\operatorname{Emb}(F)).
\end{equation}
In particular, all four selected determinants are nonzero.
\end{lemma}
\begin{proof}
Put $n_i=u(p_i-q_iv)$ and $N_i=d_n n_i$. Equation
\eqref{eq:fixedclearing} makes $N_i=(d_nu)p_i-(d_nuv)q_i$ integral.
If some $\beta_i>\gamma_i$, the matrix entry is zero. Otherwise put
$\delta_i=\gamma_i-\beta_i$. After extracting $Z^\beta$ and applying
the original factors, the entry is the embedded image of
\[
a^{jh}\prod_i\binom{\gamma_i}{\beta_i}
[t^s](1+t)^h\prod_iD_i^{e_i}(jn_i+q_iL_i(t))^{\delta_i}.
\]
The additional factor \eqref{eq:extracolumn} clears the remaining
denominators because
\begin{align*}
d_a^{Kh}a^{jh}&=d_a^{(K-j)h}(d_aa)^{jh},\\
d_n^{\gamma_i}(jn_i+q_iL_i)^{\delta_i}
 &=d_n^{\beta_i}(jN_i+d_nq_iL_i)^{\delta_i}.
\end{align*}
All exponents are nonnegative. The polynomial
$D_i(jN_i+d_nq_iL_i)$ is integral coefficientwise, and
$\delta_i\le\gamma_i\le e_i$, so its $\delta_i$th power is cleared
by $D_i^{e_i}$. Products, binomial coefficients and coefficient
extraction preserve integrality.

Thus these entries define one square matrix over $\OO_F$. Its
determinant is $Z_H$, and determinants commute with every embedding.
This proves \eqref{eq:lift}. At $\sigma_0$ its image is nonzero,
since both factors are positive and the selected minor is nonzero.
Hence $Z_H\ne0$, and injectivity of every field embedding proves the
last assertion. No new column choice is made at the other embeddings.
\end{proof}

\begin{lemma}\label{lem:cost}
With
\begin{equation}\label{eq:tau}
\tau=\frac{K\log d_a}{w_0}+\frac{\log d_n}{w_*},
\end{equation}
one has $0\le\log A_H\le MH\tau$.
\end{lemma}
\begin{proof}
Every column satisfies $h\le H/w_0$ and
$\sum_i\gamma_i\le H/w_*$. Apply these inequalities to the logarithm
of \eqref{eq:extracolumn}, then sum over $M$ columns. The logarithms
of the positive integer denominators are nonnegative. The transverse
cost has no additional factor $m$, because the column budget bounds
the sum of all exponents.
\end{proof}

The nonzero integer norm of $Z_H$ gives
\[
0\le4\log(\mathscr C_H A_H)+\sum_\sigma\log|\Delta_\sigma|.
\]
Using \eqref{eq:oldcost}, Lemma~\ref{lem:cost}, and the estimates at
two good and two remaining embeddings, then dividing by $2MH>0$,
gives
\begin{align}
-2(1-b_H)-2(E_0+\tau)-U(R)-r_c(H)
\le{}&V+r_a(H)+\frac{\log2}{w_*}\notag\\
&+\max\{-C_H,-\nu(A(1-\eta)-b_H)\}.\label{eq:combined}
\end{align}
The added normalized arithmetic cost is at most $2\tau$.
It is counted only here; the analytic bounds still concern the
unscaled determinants.

\section{Absorption and the contradiction}\label{sec:absorption}
Set
\[
E=E_0+\frac{K\log d_a}{w_0},\qquad
\widehat R=2d_n^2(R+2).
\]
For $K\ge1$, $R\ge0$, and $d_n\ge1$,
\[
2(K\widehat R+2)\ge 2d_n^2\,2(KR+2).
\]
Indeed the difference is $8d_n^2(K-1)+4\ge0$. Taking logarithms
and dividing by $w_*>0$ proves
\begin{equation}\label{eq:absorb}
U(R)+\frac{\log2+2\log d_n}{w_*}\le U(\widehat R).
\end{equation}
This enlargement pays the transverse denominator cost and the
good-embedding multiplier simultaneously.

We now justify the required parameter choices with these fixed costs.
First choose shape parameters with \eqref{eq:gap}; for example take
rational $t<1$ close enough to one that $\nu t>2(1+t)$, put
$\theta=t^2$, choose rational $\theta<A<B<t$ close to $t$, then
choose $C\in(B/\theta,1/B)$ and a sufficiently small $\eta>0$.
Choose $0<\epsilon<g$ with $\epsilon\le1/2$, then choose
$F_0>2/\theta$ with $\nu/F_0<\epsilon/3$.

Choose $m$ sufficiently large that
\begin{equation}\label{eq:dimensionmargin}
\frac{2\Lambda F_0m+3\log2}{v_0}
+\frac{(\rho+\kappa+2\log d_a)K+\log(3/2)}{w_0}
<\epsilon/3,
\qquad C_\infty>3.
\end{equation}
This is possible since $K/w_0\le(CB)^m\to0$,
$m/v_0\to0$ because $C\theta/B>1$, and
$C_\infty\to\infty$ by \eqref{eq:collision}.
Next choose the geometric slack required for interpolation and increase
the initial approximation-weight threshold so that all remaining
inverse-weight terms in $2E+U(\widehat R)+V$ sum to less than
$\epsilon/3$. Their coefficients are fixed at this point, and
$\sum_i1/w_i\le m/w_*$. Later denominators can still be chosen
successively large enough for separation. Consequently
\begin{equation}\label{eq:margins}
2E+U(\widehat R)+V<g,
\qquad 2+2E+U(\widehat R)+V<C_\infty.
\end{equation}
All these choices, including the finite approximation tuple, precede
the increase of $H$. In particular $K\log d_a/w_0$ and
$\log d_n/w_*$ are not height remainders.

By \eqref{eq:absorb}, \eqref{eq:combined} implies
\[
-2(1-b_H)\le S_H+
\max\{-C_H,-\nu(A(1-\eta)-b_H)\},
\]
where $S_H=2E+U(\widehat R)+V+r_a(H)+r_c(H)$.
For all sufficiently large $H$, \eqref{eq:margins} yields
$S_H<g$ and $2+S_H<C_H$.
If the maximum is the approximation term, rearranging gives
\[
g+(\nu-2)(\theta-b_H)
=\nu(A(1-\eta)-b_H)-2(1-b_H)\le S_H<g,
\]
a contradiction. If it is the collision term, it gives
$C_H\le2(1-b_H)+S_H\le2+S_H$, again a contradiction.
Either argument works if the two terms are equal.

Interpolation supplies a nonzero minor at a cofinal set of heights, so
this contradiction disproves the existence of arbitrarily large bad
denominators. This proves the literal bound in
Theorem~\ref{thm:general}. A rational $x$ would have exact
representations with arbitrarily large unreduced denominators, contrary
to that bound. Thus $x$ is irrational. For each $\nu>4$, only finitely
many reduced fractions can violate it: denominators are bounded, and
the inequality bounds their numerators as well. Hence $\mu(x)\le4$.

\section{Specialization and published comparison}
Take $F=\Q(\sqrt2,\sqrt3)$, which has degree four, and put
\[
a=3/2,\qquad u=1,\qquad v=\sqrt2,\qquad
x=\log(3/2)+\sqrt2>0.
\]
The two embeddings fixing $\sqrt2$ and sending $\sqrt3$ to its
two signs are distinct. They satisfy \eqref{eq:compatible}, since
$\exp(x-\sqrt2)=3/2$, and their slopes have absolute value one.
The positive distinguished base is $3/2\ne1$. We may choose
$d_a=2$, $d_n=1$ in \eqref{eq:fixedclearing}.
Theorem~\ref{thm:general} proves Theorem~\ref{thm:pilot}.

\subsection{Height conversion}\label{sec:benchmark}
For coprime $p,q$, $q>0$, the quadratic number
$U=p/q-\sqrt2$ has primitive irreducible integer polynomial
\[
q^2X^2-2pqX+(p^2-2q^2).
\]
Primitivity follows from $\gcd(q^2,p^2-2q^2)=1$; irreducibility
follows from $\sqrt2\notin\Q$. Its coefficient height obeys
$H(U)\ge q^2$. For approximations with $|x-p/q|<1$, it also obeys
$H(U)\le C_xq^2$, where $C_x>0$ is fixed.
Marcovecchio's definition and table
\cite[pp.~147, 151, 182]{marcovecchio} imply, for every
$s>9.460812$ and all sufficiently large such heights,
\[
|x-p/q|=|\log(3/2)-U|>H(U)^{-s}
\ge C_x^{-s}q^{-2s}.
\]
Choosing $9.460812<s<\nu/2$ and absorbing the fixed constant gives
the comparison bound $\mu(x)\le18.921624$.
The factor two comes from the quadratic growth of coefficient height.
The comparison checks an applicable consequence of that result;
it does not settle the strongest current bound or priority for four.

\section{Formalization and verification status}\label{sec:formalization}
The new source is maintained separately from the inherited affine
package. Its main mathematical boundary is the entry lift and the
positive scalar for the same selected determinant. The numerical
parameter structures are reused with base-denominator parameter
$\log d_a$ and coarse bound $\widehat R$.

\paragraph{Source validation.}
The general denominator-aware endpoint and the pilot corollaries have
compiled in Lean 4.34.1 with strict warning settings. The ten new source
files are identified by a SHA-256 inventory. A transitive axiom-policy
audit passed for all seven public endpoints, with rejection controls.
Separate mathematical, specification and blind manuscript reviews
found no remaining mathematical or statement defect in this scope.

\paragraph{Independent verification.}
Six statements passed comparison against a separate challenge and
both of Comparator's kernel checks. Nanoda and a separately built
stock Lean 4.34.1 reader then accepted identical retained export bytes.
The export contains 112854 serialized declarations, including proof
bodies for all six public targets and all six comparison wrappers,
and exactly the three permitted axioms. Source and export hashes,
acceptance and rejection controls, and tool identities accompany
the package. Local execution used identified inherited caches;
a fresh hosted build and its evidence audit remain pending.
The pinned Nanoda checker accepts a known invalid-projection test
that both stock Lean checks reject. Its result is supplementary;
the stock Lean replay remains required.

The standard-definition corollary excludes
\texttt{LiouvilleWith}$(\nu,x)$ for every $\nu>4$. To absorb the
constant in that definition, apply the literal bound at
$\nu'=(4+\nu)/2$ and use the frequently occurring smaller-exponent
approximations furnished by Mathlib. Irrationality is proved separately,
so the statement does not rest on a default supremum value.

\section*{Attribution and limitations}
This work reuses the interpolation, collision, logarithmic-jet and
parameter-selection machinery in the OpenAI source and the preceding
affine adaptation \cite{openai,affine}. The new step is the
fixed-denominator arithmetic extension and its specialization above.
Code provenance, licenses and exact declaration correspondence belong
to the accompanying formalization package. The present paper does not
claim a formalization of every possible quadratic-field affine family,
an effective approximation threshold, or a new best published record.

OpenAI Codex assisted with mathematical development, Lean proof
construction, verification orchestration and exposition. Separate
model contexts perform mathematical, specification and manuscript
reviews. Such reviews are not human peer review, and kernel acceptance
does not establish historical priority.

\begin{thebibliography}{9}
\raggedright
\bibitem{affine}
R. M. Casper,
\emph{\href{https://github.com/Mattie/math/blob/cf1338bcee8fe9f08422ec36e8fdf6dc6392c677/preprints/The-irrationality-exponents-of-log-2-plus-sqrt2-and-sqrt2-log-2-are-at-most-4-October-11-2026/paper.pdf}{The irrationality exponents of $\log2+\sqrt2$ and
$\sqrt2\log2$ are at most four}}, October 11, 2026.
Companion preprint and \texttt{AffineLogarithms} formalization in the
\texttt{Mattie/math} repository.

\bibitem{openai}
OpenAI, \emph{\texttt{OAI.NumberTheory.PiExponent}}, Lean source library,
revision \path{adc7f1241b42e322a6451854ab7e4b4c146bf78a},
\url{https://github.com/openai/math}.
The rational-logarithm extensions are inherited through the
\texttt{Mattie/math} revision
\path{0cfe10002ea95c45721b93d18bffbe2c3bbbcbbd}; the
\texttt{AffineLogarithms} package records their exact provenance.

\bibitem{marcovecchio}
R. Marcovecchio, \emph{The Rhin--Viola method for $\log2$},
Acta Arithmetica \textbf{139} (2009), no.~2, 147--184.
\href{https://doi.org/10.4064/aa139-2-5}{doi:10.4064/aa139-2-5}.
\end{thebibliography}
\end{document}
